\subsection{Class-Interference GNN}
\label{sec:class_interference_gnn}

Learning a new category can affect old categories unevenly.
Uniform replay does not distinguish an old category that is highly vulnerable
from one that is largely unaffected. We therefore construct a directed class
graph to estimate the interference caused by an incremental update.

For each category $c$, let $\mathcal{M}_c$ denote the Hungarian-matched
query--target pairs whose target category is $c$. Its category-conditioned
detection loss is defined as
\begin{equation}
\mathcal{L}_c
=\frac{1}{\max(1,|\mathcal{M}_c|)}
\sum_{(q,b)\in\mathcal{M}_c}
\left[
\lambda_{\mathrm{cls}}\ell_{\mathrm{cls}}(\hat{p}_q,c)
+\lambda_1\|\hat{b}_q-b\|_1
+\lambda_{\mathrm{giou}}\ell_{\mathrm{giou}}(\hat{b}_q,b)
\right],
\label{eq:category_detection_loss}
\end{equation}
where $\hat{p}_q$ and $\hat{b}_q$ are the class prediction and predicted box
of query $q$, respectively, and $b$ is its matched ground-truth box.
The terms $\ell_{\mathrm{cls}}$ and $\ell_{\mathrm{giou}}$ denote the
classification and generalized-IoU losses, and the $\lambda$ coefficients
are the corresponding detection-loss weights.
We compute the gradient of this loss with respect to the target decoder FFN
parameters as $G_c=\partial\mathcal{L}_c/\partial\Theta_{\mathrm{FFN},t}$
and vectorize it as $g_c=\operatorname{vec}(G_c)$, where
$\Theta_{\mathrm{FFN},t}$ denotes the FFN parameters selected for adaptation
at stage $t$. The resulting gradient is compressed and normalized to obtain a
fixed-dimensional node feature $x_c$. Each category is represented by one node,
and an ordered pair of nodes represents a possible directed interference edge.

To construct supervision for the graph, we perform a small probe update using
the gradient of the newly introduced category $u$, yielding
$\Theta_{\mathrm{FFN},t,u}^{\mathrm{probe}}
=\Theta_{\mathrm{FFN},t}-\eta\,\bar{G}_u$,
where $\bar{G}_u$ is the normalized gradient reshaped to match the target
FFN parameters and $\eta$ is a small probe step size. The empirical harm
from category $u$ to an old category $c$ is measured by the decrease in its
validation AP:
\begin{equation}
H_{u\rightarrow c}
=\left[\operatorname{AP}_{c}(f_{\theta_t})
-\operatorname{AP}_{c}(f_{\theta_{t,u}^{\mathrm{probe}}})\right]_{+},
\label{eq:empirical_harm}
\end{equation}
where $\theta_t$ is the detector before the probe update and
$\theta_{t,u}^{\mathrm{probe}}$ is the detector obtained after replacing the
target FFN parameters with $\Theta_{\mathrm{FFN},t,u}^{\mathrm{probe}}$.
The notation $[x]_{+}=\max(x,0)$ keeps only positive AP degradation.

The GNN receives the class-level gradient features and performs message passing
between category nodes. At layer $l$, the representation of category $i$ is
updated as
\begin{equation}
\begin{aligned}
h_i^{(l+1)}=h_i^{(l)}+\Phi_l\Bigl(
h_i^{(l)},\sum_{j\neq i}a_{j\rightarrow i}^{(l)}
\Psi_l\bigl([h_j^{(l)},h_i^{(l)},h_j^{(l)}-h_i^{(l)}]\bigr)
\Bigr),
\end{aligned}
\label{eq:gnn_message_passing}
\end{equation}
where $h_i^{(0)}$ is obtained from $x_i$ and $a_{j\rightarrow i}^{(l)}$
is a learned normalized coefficient for the directed edge from category $j$
to category $i$. The edge representation contains the source node, the target
node, and their difference, allowing the graph to model asymmetric interference.

After the final message-passing layer, the GNN predicts the interference
probability from source category $i$ to target category $j$:
\begin{equation}
p_{i\rightarrow j}
=\sigma\left(\Psi_{\mathrm{edge}}\left(
[h_i^{(L)},h_j^{(L)},h_i^{(L)}-h_j^{(L)}]
\right)\right).
\label{eq:edge_prediction}
\end{equation}

The GNN is trained using interference artifacts collected from previously
completed incremental stages. Let $\mathcal{E}_{\mathrm{hist}}$ denote the
measured directed edges in these artifacts. The measured harm values in
Eq.~\ref{eq:empirical_harm} are normalized within each source category as
$\widetilde{H}_{u\rightarrow c}
=H_{u\rightarrow c}/\max(\epsilon,\max_v H_{u\rightarrow v})$,
where the maximum is taken over measured targets in the same artifact and
$\epsilon>0$ prevents division by zero. Let $r_{i\rightarrow j}$ denote
the edge logit before the sigmoid in Eq.~\ref{eq:edge_prediction}, and let
$\mathcal{Q}_{\mathrm{hist}}$ contain within-artifact triplets
$(u,c^+,c^-)$ satisfying $H_{u\rightarrow c^+}>H_{u\rightarrow c^-}$.
The GNN is optimized with
\begin{equation}
\begin{aligned}
\mathcal{L}_{\mathrm{GNN}}
={}&\frac{1}{|\mathcal{E}_{\mathrm{hist}}|}
\sum_{(u,c)\in\mathcal{E}_{\mathrm{hist}}}
\operatorname{SmoothL1}\left(p_{u\rightarrow c},
\widetilde{H}_{u\rightarrow c}\right)\\
&+\frac{\lambda_{\mathrm{rank}}}{\max(1,|\mathcal{Q}_{\mathrm{hist}}|)}
\sum_{(u,c^+,c^-)\in\mathcal{Q}_{\mathrm{hist}}}
\left[m_{\mathrm{rank}}-r_{u\rightarrow c^+}
+r_{u\rightarrow c^-}\right]_{+},
\end{aligned}
\label{eq:gnn_loss}
\end{equation}
where the regression term fits the normalized harm values, while the pairwise
ranking term encourages categories with larger measured AP degradation to
receive larger edge scores. The hyperparameters $\lambda_{\mathrm{rank}}$
and $m_{\mathrm{rank}}$ control the ranking weight and margin, respectively.
The GNN parameters are updated by minimizing $\mathcal{L}_{\mathrm{GNN}}$.
The current stage is not used to train its own neighborhood predictor; its
measured interference is stored for later graph training.

At task $T_{t+1}$, the GNN predicts an interference score $p_{u\rightarrow c}$
for every ordered pair $(u,c)\in K_{t+1}\times\mathcal{K}_t$.
Because multiple new categories are introduced simultaneously, we aggregate
their predicted effects on each old category as
$s_{t+1}(c)=\max_{u\in K_{t+1}}p_{u\rightarrow c}$ for
$c\in\mathcal{K}_t$. The maximum operator treats an old category as vulnerable
if it is strongly affected by any category introduced at the current task.
The old categories with the largest aggregated interference scores form the
task-level neighborhood:
\begin{equation}
\mathcal{N}_{t+1}
=\operatorname{TopK}_{c\in\mathcal{K}_t}s_{t+1}(c).
\label{eq:interference_neighborhood}
\end{equation}
The selected $\mathcal{N}_{t+1}$ is used to determine how historical replay
samples are allocated and which old categories receive local protection during
the incremental update.

% Preamble: \usepackage{algorithm,algpseudocode}
\begin{algorithm}[t]
\caption{Training the Class-Interference GNN}
\label{alg:gnn_training}
\begin{algorithmic}[1]
\renewcommand{\algorithmicrequire}{\textbf{Input:}}
\renewcommand{\algorithmicensure}{\textbf{Output:}}
\Require Completed historical stages $\mathcal{S}_{\mathrm{hist}}$ with
pre-update detectors, class data, and validation data; GNN parameters $\phi$;
probe step size $\eta$; GNN learning rate $\eta_g$; training epochs $E$
\Ensure Trained class-interference GNN $G_\phi$
\State Initialize the historical artifact set $\mathcal{A}\gets\emptyset$
\For{each stage $s\in\mathcal{S}_{\mathrm{hist}}$}
    \State Load its pre-update detector $f_{\theta_s}$
    \For{each category $c$ observed at stage $s$}
        \State Compute $\mathcal{L}_c$ using Eq.~\ref{eq:category_detection_loss}
        \State Compute $g_c\gets\operatorname{vec}(\nabla_{\Theta_{\mathrm{FFN},s}}\mathcal{L}_c)$
        \State Compress and normalize $g_c$ to obtain node feature $x_c$
    \EndFor
    \For{each newly introduced category $u$ at stage $s$}
        \State Apply a temporary probe:
        $\Theta_{\mathrm{FFN},s,u}^{\mathrm{probe}}
        \gets\Theta_{\mathrm{FFN},s}-\eta\bar{G}_u$
        \State Measure $H_{u\rightarrow c}$ for each old category $c$
        using Eq.~\ref{eq:empirical_harm}
        \State Restore the pre-update detector $f_{\theta_s}$
    \EndFor
    \State Store the node features and measured harm in $\mathcal{A}$
\EndFor
\For{$e=1,\ldots,E$}
    \For{each historical artifact in $\mathcal{A}$}
        \State Predict $p_{u\rightarrow c}$ using
        Eqs.~\ref{eq:gnn_message_passing} and~\ref{eq:edge_prediction}
        \State Compute $\mathcal{L}_{\mathrm{GNN}}$ using Eq.~\ref{eq:gnn_loss}
        \State Update $\phi\gets\phi-\eta_g\nabla_\phi\mathcal{L}_{\mathrm{GNN}}$
    \EndFor
\EndFor
\State Freeze $G_\phi$ for neighborhood prediction at the current stage
\State \Return $G_\phi$
\end{algorithmic}
\end{algorithm}

\subsection{Graph-Guided Detector Adaptation}
\label{sec:graph_guided_detector_adaptation}

Updating the complete detector during an incremental stage can alter
representations that are unrelated to the new category. We therefore use the
predicted neighborhood $\mathcal{N}_u$ to determine which detector parameters
are exposed to the incremental update and which parameter directions remain
protected.

Let $\Theta_t$ denote the detector parameters after stage $T_t$ and let
$\mathcal{P}(\mathcal{N}_u)$ be the parameter subset selected by the graph for
the neighborhood of $u$. The graph-guided update is written as
\begin{equation}
\widetilde{\Theta}_t=\Theta_t+\Delta\Theta_{\mathcal{N}_u},
\qquad
\operatorname{supp}(\Delta\Theta_{\mathcal{N}_u})
\subseteq\mathcal{P}(\mathcal{N}_u),
\label{eq:graph_guided_update}
\end{equation}
where $\Delta\Theta_{\mathcal{N}_u}$ is the update induced by the current
stage objective and $\operatorname{supp}(\cdot)$ denotes the parameter entries
that receive nonzero updates. The selection is made at the detector level and
does not require a separate parameter bank for each category.

Parameters outside $\mathcal{P}(\mathcal{N}_u)$ are protected during the
incremental step, while parameters selected for the neighborhood receive
gradients from the new-category and replay supervision. This class-conditioned
parameter routing prevents the update from treating every old category as
equally vulnerable.

The selected detector parameters are optimized using the labeled data for the
new category, historical replay data from the predicted neighborhood, and
high-confidence predictions of the frozen detector. Let $\mathcal{D}_u$ denote
the labeled data for the new category. Teacher completion produces
$\widetilde{\mathcal{D}}_u
=\mathcal{D}_u\cup\widehat{\mathcal{D}}_{\mathrm{old}}(\mathcal{D}_u)$,
where $\widehat{\mathcal{D}}_{\mathrm{old}}$ contains high-confidence
predictions of old objects generated by the frozen detector. This operation
prevents old foreground objects that are absent from the current annotations
from being treated as background.

Replay is sampled from historical images containing categories in
$\mathcal{N}_u$. The replay budget is fixed across methods, while the graph
determines how that budget is distributed over old categories. This gives the
selected parameters access to examples that are most relevant to the predicted
interference pattern.

After the incremental optimization described in
Section~\ref{sec:graph_conditioned_objective}, the selected update is applied
to obtain the detector for the next stage as
$\Theta_{t+1}=\Theta_t+\Delta\Theta_{\mathcal{N}_u}$.
The resulting detector is used as the base model for the next incremental
stage, while the graph is updated with the observed class-level interference.

\subsection{Graph-Conditioned Continual Objective}
\label{sec:graph_conditioned_objective}

The incremental detector is optimized with the standard DETR detection loss
and two graph-conditioned regularization terms. The detection loss is
computed on the teacher-completed current data and the replay set selected
from the predicted neighborhood as
$\mathcal{L}_{\mathrm{det}}
=\mathcal{L}_{\mathrm{DET}}(\widetilde{\mathcal{D}}_u
\cup\mathcal{R}(\mathcal{N}_u))$,
where $\mathcal{R}(\mathcal{N}_u)$ denotes the historical replay set restricted
to categories in $\mathcal{N}_u$.

The first additional term improves discrimination between the new category and
the old categories in its predicted neighborhood. For a matched query $q$
with target category $y$, the correct category logit is required to exceed the
strongest competing local category by a margin $m$:
\begin{equation}
\mathcal{L}_{\mathrm{local}}
=\frac{1}{|\mathcal{M}|}
\sum_{(q,y)\in\mathcal{M}}
\left[m-\left(z_{q,y}
-\max_{\substack{c\in\mathcal{N}_u\cup\{u\}\\c\neq y}}
z_{q,c}\right)\right]_{+},
\label{eq:local_margin}
\end{equation}
where $\mathcal{M}$ contains matched queries whose target categories belong to
$\mathcal{N}_u\cup\{u\}$. This term focuses on local class competition and
does not impose an artificial margin between every pair of categories.

The graph-guided update may still overlap with parameter directions that are
important for categories outside $\mathcal{N}_u$. We therefore construct an
orthonormal basis $U_{\mathrm{off}}$ from the class-conditioned gradients of
non-neighborhood categories. Let
$\delta=\operatorname{vec}(\Delta\Theta_{\mathcal{N}_u})$
denote the vectorized graph-guided update, where updates from all selected
parameter groups are concatenated when necessary. The off-neighborhood
protection term is defined as
\begin{equation}
\mathcal{L}_{\mathrm{off}}
=\left\|U_{\mathrm{off}}^{\top}\delta\right\|_2^2.
\label{eq:off_neighborhood}
\end{equation}

The complete incremental objective is
\begin{equation}
\mathcal{L}_{\mathrm{total}}
=\mathcal{L}_{\mathrm{det}}
+\lambda_{\mathrm{local}}\mathcal{L}_{\mathrm{local}}
+\lambda_{\mathrm{off}}\mathcal{L}_{\mathrm{off}}.
\label{eq:total_objective}
\end{equation}

The four protection mechanisms are integrated into this objective and its data
construction. First, graph-guided parameter routing confines learning to the
selected detector parameters. Second, teacher completion preserves old
foreground objects that are not annotated in the current incremental data.
Third, neighborhood replay revisits historical examples from categories
predicted to be at risk. Finally, off-neighborhood projection restricts the
update from overlapping with directions associated with categories outside
the predicted neighborhood.

The local margin term is not an independent forgetting protection. It improves
the separation between the new category and the old categories that are most
likely to compete with it. After the objective is optimized, the graph-guided
update is applied as described in
Section~\ref{sec:graph_guided_detector_adaptation}, and the resulting detector
is passed to the next incremental stage.
